\section{The ideal tetrahedron: an explicit volume integral}
\label{sec:tetrahedron}

\subsection{Lambda lengths and the ambient cross-ratio}

Take the four labeled ideal vertices to be $0,1,\infty,z$, where
$z=x+iy\in\CC\setminus\{0,1\}$. For $y\ne0$ their convex hull is a genuine
ideal tetrahedron. The two half-planes describe its two orientation types.

For a decorated ideal triangulation, set $\lambda_e=e^{\delta_e/2}$, where
$\delta_e$ is the signed distance between the chosen horocycles along the
edge. An overall fixed multiplier in the definition of $\lambda$ does not
change its logarithmic differential. With cyclically ordered edge lengths
$a,b,c$ on each oriented triangle, Penner's formula in convention
\eqref{eq:fn} is
\begin{equation}
 \omega_{\WP}=-2\sum_{\triangle}
 \bigl(\dd\log a\wedge\dd\log b+
       \dd\log b\wedge\dd\log c+
       \dd\log c\wedge\dd\log a\bigr).
 \label{eq:penner}
\end{equation}
This is a form on decorated Teichm\"uller space that descends after forgetting
the horocycles~\cite{Penner}; its magnitude agrees with the cuspidal limit
of Mondello's arc-coordinate formula~\cite{MondelloBoundary}.

The four decorations can be chosen so that
\begin{equation}
 \lambda_{0\infty}=\lambda_{1\infty}=\lambda_{z\infty}
 =\lambda_{01}=1,
 \qquad \lambda_{0z}=r=|z|,\quad\lambda_{1z}=s=|1-z|.
 \label{eq:lambdagauge}
\end{equation}
For example, use equal Euclidean diameters for the three finite horospheres
in the upper-half-space model and adjust their common scale together with
the horosphere at infinity. The usual formula for the signed distance then
gives~\eqref{eq:lambdagauge}.
One may take diameter $1$ for each finite horosphere and height $1$ at
infinity. The finite-edge formula is then
$\lambda_{ij}=|z_i-z_j|$, and each edge to infinity has lambda length $1$.
Overlapping horospheres cause no problem: the distances are signed, and
smaller embedded decorations can first be chosen and then rescaled.

Every face containing $\infty$ has at most one varying lambda length, so its
contribution in~\eqref{eq:penner} vanishes in this section of the decoration
bundle. The remaining face has lengths $1,r,s$. Consequently the positive
WP measure is
\begin{equation}
 \boxed{
 \dV=2\abs{\dd\log r\wedge\dd\log s}
      =\frac{2|y|}{|z|^2|1-z|^2}\,\dd x\,\dd y.}
 \label{eq:tetdensity}
\end{equation}
The Jacobian calculation is elementary:
\[
 \dd\log r=\frac{x\,\dd x+y\,\dd y}{r^2},\qquad
 \dd\log s=\frac{(x-1)\,\dd x+y\,\dd y}{s^2}.
\]
Their wedge is $y\,\dd x\wedge\dd y/(r^2s^2)$.

\subsection{Integration in angle coordinates}

For $y>0$, let $A,B,C$ be the interior angles of the Euclidean triangle
with vertices $0,1,z$. They satisfy
\[
 A>0,\quad B>0,\quad C>0,\quad A+B+C=\pi.
\]
They are also the three interior dihedral angles of the ideal tetrahedron,
each repeated on the opposite edge. Direct differentiation gives
\[
 \abs{\dd A\wedge\dd B}
   =\frac{y}{r^2s^2}\,\dd x\,\dd y,
 \qquad\text{hence}\qquad
 \dV=2\,\dd A\,\dd B.
\]
One orientation sheet therefore contributes
\begin{equation}
 \int_{\substack{A>0,\ B>0\\A+B<\pi}}2\,\dd A\,\dd B=\pi^2.
 \label{eq:onecell}
\end{equation}
Including the other sheet gives the promised direct calculation:
\begin{equation}
 \boxed{V_{0,4}=2\pi^2.}
 \label{eq:v04direct}
\end{equation}
The real configurations form the flat locus and contribute zero two-dimensional
WP measure.

There is also a useful check on~\eqref{eq:polyhedralpartition}. In a
tetrahedral angle cell, the four perimeter equations force opposite exterior
angles to be equal, with three values $a,b,c$ satisfying $a+b+c=2\pi$.
To see the disintegration factor explicitly, order the edges as
$(01,02,03,12,13,23)$, keep $a=\theta_{01}$ and $b=\theta_{02}$,
and eliminate the last four variables. The corresponding perimeter minor is
\[
 \begin{pmatrix}
 1&0&0&0\\0&1&1&0\\0&1&0&1\\1&0&1&1
 \end{pmatrix},\qquad \det=-2.
\]
It gives $(\theta_{03},\theta_{12},\theta_{13},\theta_{23})
=(2\pi-a-b,2\pi-a-b,b,a)$ and the measure
$\frac12|\dd a\,\dd b|$. Up to permuting the three angle names,
$a=\pi-A$ and $b=\pi-B$, so this is
$\frac12\dd A\,\dd B$. There are two oriented cells, each with trivial
label-preserving automorphism group. The factor $2^{4-2}=4$ in
\eqref{eq:polyhedralpartition} converts their total map volume $\pi^2/2$
to $2\pi^2$.

\subsection{What happens at the punctures of the parameter plane?}

Finite volume requires integrability, not pointwise decay of a coordinate
density. Near $z=0$, write $z=\rho e^{i\vartheta}$. Then
\[
 \frac{2|y|}{|z|^2|1-z|^2}
 =\frac{2|\sin\vartheta|}{\rho}\,(1+O(\rho)),
 \qquad \dd x\,\dd y=\rho\,\dd\rho\,\dd\vartheta.
\]
The singularity is integrable. The same is true at $1$. At infinity the
density is $O(|z|^{-3})$, again integrable against Euclidean area. Thus the
requested decay is best formulated as a condition on the measure, not on
its density without its coordinate area element.

\begin{remark}[The coordinate is not conformal]
Formula~\eqref{eq:tetdensity} is the pullback to the ambient ideal-vertex
coordinate. It is not a formula for WP measure in the intrinsic conformal
cross-ratio. In particular, its vanishing on the real locus does not mean
that the WP form degenerates at an ordinary interior point of moduli space.
The polyhedral parametrization is not a smooth coordinate across that locus.
The Delaunay literature also distinguishes the embedding coordinate from
uniformization~\cite[Section 3.3.1]{CDE}.
\end{remark}

Nor do two ordinary closed angle triangles describe the full topology of
the compactification without additional data. As $y\downarrow0$ with
$0<x<1$, two of the triangle angles tend to zero and the third to $\pi$,
independently of $x$. A whole family of flat quadrilaterals approaches one
angle-simplex vertex. Retaining that modulus requires a refinement of the
boundary, such as a real blow-up. This has no effect on the integral over
the open cells.

